% =============================================================
% main.tex
% Thesis template - Program Studi S2 Data Science dan Kecerdasan Buatan
% Fakultas Ilmu Komputer dan Teknologi Informasi (Fasilkom-TI)
% Universitas Sumatera Utara (USU)
%
% Based on "Pedoman Penulisan Tesis" Fasilkom-TI USU
% and the "Overleaf Online LaTeX Editor untuk Penulisan Tesis" guide.
% =============================================================
\documentclass[12pt,a4paper]{report}
% =============================================================
% preamble.tex
% Formatting rules implementing "Pedoman Penulisan Tesis"
% Fasilkom-TI, Universitas Sumatera Utara
% (adapted for Program Studi S2 Data Science dan Kecerdasan Buatan)
% =============================================================

% ---- Page geometry -------------------------------------------------
% Body pages : top 30mm, left 38mm, right 25mm, bottom 25mm
% Chapter-opening pages need top 50mm; handled in \chapter redefinition below.
\usepackage[a4paper,top=30mm,left=38mm,right=25mm,bottom=25mm,headsep=8mm,headheight=15pt]{geometry}

% ---- Fonts -----------------------------------------------------------
\usepackage[T1]{fontenc}
\usepackage[utf8]{inputenc}
\usepackage{times}            % Times New Roman substitute (body text 12pt)

% Localized names for the automatically generated lists (kept manual,
% rather than via babel, so the template compiles on any TeX
% distribution without requiring the Indonesian hyphenation patterns).
\renewcommand{\contentsname}{Daftar Isi}
\renewcommand{\listtablename}{Daftar Tabel}
\renewcommand{\listfigurename}{Daftar Gambar}
\renewcommand{\bibname}{Daftar Pustaka}
\renewcommand{\appendixname}{Lampiran}

% ---- Spacing -----------------------------------------------------------
\usepackage{setspace}
\onehalfspacing                % 1.5 line spacing for main text
% Tables, bibliography and other lists use \singlespacing locally.

% ---- Paragraphs --------------------------------------------------------
% First paragraph after any heading is not indented (LaTeX default behaviour).
% Subsequent paragraphs are indented 6 characters (~1.27 cm at 12pt Times).
\setlength{\parindent}{1.27cm}
\setlength{\parskip}{0pt}

% ---- Useful packages ----------------------------------------------------
\usepackage{graphicx}
\usepackage{amsmath,amssymb,amsfonts}
\usepackage{booktabs}
\usepackage{multirow}
\usepackage{array}
\usepackage{caption}
\usepackage{float}
\usepackage{longtable}
\usepackage{listings}
\usepackage{xcolor}
\usepackage{tikz}
\usepackage[hidelinks]{hyperref}

\hypersetup{
    colorlinks   = true,
    linkcolor    = black,
    citecolor    = black,
    urlcolor     = blue,
    pdftitle     = {Judul Tesis},
    pdfauthor    = {Nama Mahasiswa},
    pdfsubject   = {Tesis Magister Data Science dan Kecerdasan Buatan},
}

% ---- Citations: "Sistem Pengarang-Tahun" (author-year), alphabetical, hanging indent
\usepackage[authoryear,round]{natbib}
\bibliographystyle{apalike}
\setlength{\bibsep}{0pt}
\setlength{\bibhang}{1.27cm}   % hanging indent for reference list

% ---- Captions: "Tabel"/"Gambar" + chapter.number, centered above table / below figure
\usepackage{chngcntr}
\captionsetup[table]{name=Tabel,labelsep=period,justification=centering,singlelinecheck=false}
\captionsetup[figure]{name=Gambar,labelsep=period,justification=centering}
\counterwithin{table}{chapter}
\counterwithin{figure}{chapter}

% =============================================================
% Heading formatting
% =============================================================
\usepackage{titlesec}

% --- Chapter (BAB) ------------------------------------------------------
% Two centered bold lines: "BAB <n>" then "<JUDUL BAB>" (uppercase),
% with 2-line spacing between them and 4-line spacing before the first
% sub-heading, as illustrated in Lampiran 3 of the Pedoman.
\titleformat{\chapter}[display]
  {\normalfont\bfseries\centering}
  {\Large BAB \thechapter}
  {12pt}
  {\Large\MakeUppercase}
% "before" length below (20mm) extends the 30mm body top margin to the
% required 50mm on chapter-opening pages; "after" (28pt) leaves roughly
% 4 line-spacings before the first sub-heading, per Lampiran 3.
\titlespacing*{\chapter}{0pt}{20mm}{28pt}

% Unnumbered chapters (\chapter*, used internally by \tableofcontents,
% \listoffigures, \listoftables and \bibliography) keep the normal 30mm
% top margin -- only numbered BAB pages get the extra 20mm.
\titleformat{name=\chapter,numberless}[display]
  {\normalfont\bfseries\centering}
  {}
  {0pt}
  {\Large\MakeUppercase}
\titlespacing*{name=\chapter,numberless}{0pt}{0pt}{28pt}

% --- Section = sub-bab level 1 -------------------------------------------
% Bold, Title Case, flush left at the chapter margin, numbered "N.M"
\titleformat{\section}
  {\normalfont\bfseries}
  {\thesection}
  {0.5em}
  {}
\titlespacing*{\section}{0pt}{18pt}{8pt}

% --- Subsection = sub-bab level 2 ----------------------------------------
% Italic, sentence case, numbered "N.M.K"
\titleformat{\subsection}
  {\normalfont\itshape}
  {\thesubsection}
  {0.5em}
  {}
\titlespacing*{\subsection}{0pt}{14pt}{6pt}

% --- Subsubsection (beyond the two mandated levels; used sparingly) -------
\titleformat{\subsubsection}
  {\normalfont\itshape}
  {\thesubsubsection}
  {0.5em}
  {}
\titlespacing*{\subsubsection}{0pt}{12pt}{4pt}

% =============================================================
% Page numbering: Times New Roman 12pt, top-right corner,
% 15 mm from the top edge, right-aligned 25 mm from the right edge.
% =============================================================
\usepackage{fancyhdr}
\pagestyle{fancy}
\fancyhf{}
\fancyhead[R]{\thepage}
\renewcommand{\headrulewidth}{0pt}
\renewcommand{\footrulewidth}{0pt}

% Chapter-opening pages ("plain" style) show no header/footer content,
% but the page is still counted in the numbering sequence.
\fancypagestyle{plain}{%
  \fancyhf{}
  \renewcommand{\headrulewidth}{0pt}
  \renewcommand{\footrulewidth}{0pt}
}

% Front-matter "title-style" pages (cover, approval, originality, etc.)
% also print no visible page number, while the roman counter keeps counting.
\newcommand{\unnumberedpage}{\thispagestyle{empty}}

% ---- Widow/orphan control (two-line minimum per the "no orphan / no widow" rule)
\clubpenalty = 10000
\widowpenalty = 10000
\displaywidowpenalty = 10000

% ---- Table of contents depth: show up to subsection ----------------------
\setcounter{tocdepth}{2}
\setcounter{secnumdepth}{3}

% ---- Faculty logo (50 x 50 mm, per Lampiran 1/4) --------------------------
% If figures/logo_usu.png is ever missing/renamed, swap this back to the
% \fbox placeholder below so the template still compiles.
\newcommand{\ThesisLogo}{%
  \includegraphics[width=50mm,height=50mm,keepaspectratio]{figures/logo_usu.png}%
  % \fbox{\parbox[c][50mm][c]{50mm}{\centering\footnotesize LOGO\\FAKULTAS ILMU KOMPUTER\\DAN TEKNOLOGI INFORMASI\\USU}}%
}


\begin{document}

% ---------------------------------------------------------------
% BAGIAN MUKA (front matter) -- roman page numbers, counted from
% the title page even though most of these pages print no number.
% ---------------------------------------------------------------
\pagenumbering{roman}
% =============================================================
% Kulit Depan (outer cover), based on Lampiran 1 of the Pedoman.
% All text: Times New Roman 14pt, uppercase, black, centered.
% Top/bottom margin on the cover: 60 mm from the paper edge.
% =============================================================
% NB: intentionally not using the \titlepage environment here -- it
% resets \c@page to 1 on every use, which would collide with the
% roman-numbered title page that follows. \clearpage + \thispagestyle{empty}
% gives the same unnumbered full page without touching the counter.
\clearpage
\thispagestyle{empty}
% The 60mm figures below are measured from the physical paper edge
% (per Lampiran 1), not from the 30mm body margin, so this page
% temporarily uses a zero margin while it is being typeset.
\newgeometry{margin=0mm}
\begin{center}
\vspace*{60mm minus 10mm}

{\fontsize{14}{18}\selectfont\bfseries JUDUL TESIS UNTUK PROGRAM MAGISTER (S2)\\
DATA SCIENCE DAN KECERDASAN BUATAN\par}

\vspace{20mm}

{\fontsize{14}{18}\selectfont\bfseries TESIS\par}

\vspace{20mm}

{\fontsize{14}{18}\selectfont\bfseries NAMA MAHASISWA\\
NIM 000000000\par}

\vfill

\ThesisLogo

\vfill

{\fontsize{14}{18}\selectfont\bfseries
PROGRAM STUDI S2 DATA SCIENCE DAN KECERDASAN BUATAN\\
FAKULTAS ILMU KOMPUTER DAN TEKNOLOGI INFORMASI\\
UNIVERSITAS SUMATERA UTARA\\
MEDAN\\
\the\year\par}

\vspace*{60mm minus 10mm}
\end{center}
\restoregeometry
\clearpage

% =============================================================
% Halaman Judul (title page), based on Lampiran 4 of the Pedoman.
% Identical to the cover, plus the "diajukan untuk..." purpose line.
% =============================================================
% See the note in cover.tex: the \titlepage environment is avoided
% because it resets \c@page, which must stay on the roman sequence here.
\clearpage
\thispagestyle{empty}
% See the note in cover.tex: 60mm here is measured from the paper edge.
\newgeometry{margin=0mm}
\begin{center}
\vspace*{60mm minus 10mm}

{\fontsize{14}{18}\selectfont\bfseries JUDUL TESIS UNTUK PROGRAM MAGISTER (S2)\\
DATA SCIENCE DAN KECERDASAN BUATAN\par}

\vspace{14mm}

{\fontsize{14}{18}\selectfont\bfseries TESIS\par}

\vspace{6mm}

{\fontsize{12}{16}\selectfont
Diajukan untuk melengkapi tugas dan memenuhi syarat memperoleh ijazah\\
Magister Data Science dan Kecerdasan Buatan\par}

\vspace{14mm}

{\fontsize{14}{18}\selectfont\bfseries NAMA MAHASISWA\\
NIM 000000000\par}

\vfill

\ThesisLogo

\vfill

{\fontsize{14}{18}\selectfont\bfseries
PROGRAM STUDI S2 DATA SCIENCE DAN KECERDASAN BUATAN\\
FAKULTAS ILMU KOMPUTER DAN TEKNOLOGI INFORMASI\\
UNIVERSITAS SUMATERA UTARA\\
MEDAN\\
\the\year\par}

\vspace*{60mm minus 10mm}
\end{center}
\restoregeometry
\clearpage

% =============================================================
% Lembar Persetujuan, based on Lampiran 5 of the Pedoman.
% =============================================================
\thispagestyle{empty}
\begin{center}
{\bfseries PERSETUJUAN\par}
\vspace{18mm} % approx. 3 x 1.5 line spacing
\end{center}

\noindent
\begin{tabular}{@{}p{4.6cm}@{}l@{\;}p{8.8cm}@{}}
Judul          & : & Judul Tesis untuk Program Magister (S2) Data Science dan Kecerdasan Buatan \\[4pt]
Kategori       & : & Tesis \\[4pt]
Nama           & : & Nama Mahasiswa \\[4pt]
Nomor Induk Mahasiswa & : & 000000000 \\[4pt]
Program Studi  & : & S2 Data Science dan Kecerdasan Buatan \\[4pt]
Fakultas       & : & Ilmu Komputer dan Teknologi Informasi\\
               &   & Universitas Sumatera Utara \\[4pt]
Komisi Pembimbing & : & \\
\end{tabular}

\vspace{12mm}

\noindent
\begin{tabular}{@{}p{7.6cm}p{7.6cm}@{}}
Pembimbing 2 & Pembimbing 1 \\[14mm]
\hrulefill & \hrulefill \\
Nama Lengkap dan Gelar Akademik & Nama Lengkap dan Gelar Akademik \\
\end{tabular}

\vspace{10mm}

\noindent Diketahui/disetujui oleh\\
Program Studi S2 Data Science dan Kecerdasan Buatan\\
Ketua,

\vspace{12mm}

\noindent\hrulefill\\
Nama Lengkap dan Gelar Akademik\\
Nomor Induk Pegawai
\clearpage

% =============================================================
% Lembar Pernyataan Orisinalitas, based on Lampiran 6 of the Pedoman.
% =============================================================
\thispagestyle{empty}
\begin{center}
{\bfseries PERNYATAAN\par}
\vspace{18mm}
\textit{Judul Tesis untuk Program Magister (S2) Data Science dan Kecerdasan Buatan}\par
\vspace{4mm}
TESIS
\end{center}

\vspace{10mm}
\noindent Saya mengakui bahwa tesis ini adalah hasil karya saya sendiri, kecuali beberapa kutipan dan ringkasan yang masing-masing telah disebutkan sumbernya.

\vspace{18mm}
\noindent Medan, Hari Bulan Tahun

\vspace{24mm}
\noindent Nama Mahasiswa\\
Nomor Induk Mahasiswa
\clearpage

% =============================================================
% Pernyataan Persetujuan Publikasi, based on Lampiran 7 of the Pedoman.
% =============================================================
\thispagestyle{empty}
\begin{center}
{\bfseries PERNYATAAN PERSETUJUAN PUBLIKASI\\
KARYA ILMIAH UNTUK KEPENTINGAN AKADEMIS\par}
\vspace{18mm}
\end{center}

\noindent Sebagai sivitas akademika Universitas Sumatera Utara, saya yang bertanda tangan di bawah ini:

\vspace{6mm}
\noindent
\begin{tabular}{@{}p{3.5cm}@{}l@{\;}l@{}}
Nama            & : & Nama Mahasiswa \\
NIM             & : & 000000000 \\
Program Studi   & : & S2 Data Science dan Kecerdasan Buatan \\
Jenis Karya Ilmiah & : & Tesis \\
\end{tabular}

\vspace{6mm}
\noindent Demi pengembangan ilmu pengetahuan, menyetujui untuk memberikan kepada Universitas Sumatera Utara Hak Bebas Royalti Non-Eksklusif (\textit{Non-Exclusive Royalty Free Right}) atas tesis saya yang berjudul:

\begin{center}
\textit{Judul Tesis untuk Program Magister (S2) Data Science dan Kecerdasan Buatan}
\end{center}

\noindent Beserta perangkat yang ada (jika diperlukan). Dengan Hak Bebas Royalti Non-Eksklusif ini, Universitas Sumatera Utara berhak menyimpan, mengalih media, memformat, mengelola dalam bentuk \textit{database}, merawat dan mempublikasikan tesis saya tanpa meminta izin dari saya selama tetap mencantumkan nama saya sebagai penulis dan sebagai pemegang dan/atau sebagai pemilik hak cipta.

\vspace{6mm}
\noindent Demikian pernyataan ini saya buat dengan sebenarnya.

\vspace{18mm}
\noindent Medan, hh-bulan-tahun

\vspace{18mm}
\noindent Nama Jelas\\
NIM
\clearpage

% =============================================================
% Lembar Panitia Penguji Tesis, based on Lampiran 8 of the Pedoman.
% =============================================================
\thispagestyle{empty}

\noindent Telah diuji pada\\
Tanggal: hh bulan tahun
\vspace{2mm}
\noindent\hrulefill

\vspace{24mm}
\noindent
\begin{tabular}{@{}p{3.2cm}l@{}}
\multicolumn{2}{@{}l}{PANITIA PENGUJI TESIS} \\[6pt]
Ketua   & : Nama Lengkap dan Gelar Akademik \\
Anggota & : 1.~Nama Lengkap dan Gelar Akademik \\
        & \phantom{1.}2.~Nama Lengkap dan Gelar Akademik \\
        & \phantom{1.}3.~Nama Lengkap dan Gelar Akademik \\
        & \phantom{1.}4.~Nama Lengkap dan Gelar Akademik \\
\end{tabular}
\clearpage

% =============================================================
% Riwayat Hidup, based on Lampiran 9 of the Pedoman.
% =============================================================
\thispagestyle{empty}
\begin{center}
{\bfseries RIWAYAT HIDUP\par}
\vspace{10mm}
\end{center}

\noindent\textbf{DATA PRIBADI}

\vspace{4mm}
\noindent
\begin{tabular}{@{}p{5.5cm}l@{}}
Nama Lengkap (berikut gelar) & : \\
Tempat dan Tanggal Lahir     & : \\
Alamat Rumah                 & : \\
Telepon/Faks/HP              & : \\
E-mail                       & : \\
Instansi Tempat Bekerja      & : \\
Alamat Kantor                & : \\
\end{tabular}

\vspace{10mm}
\noindent\textbf{DATA PENDIDIKAN}

\vspace{4mm}
\noindent
\begin{tabular}{@{}p{2cm}l l@{}}
SD   & : & TAMAT: tahun \\
SLTP & : & TAMAT: tahun \\
SLTA & : & TAMAT: tahun \\
S1   & : & TAMAT: tahun \\
S2   & : Data Science dan Kecerdasan Buatan, USU & TAMAT: tahun \\
\end{tabular}
\clearpage

% =============================================================
% Ucapan Terima Kasih (Bab 3.9 of the Pedoman).
% Keep this concise; thank persons/institutions that helped
% materially or non-materially.
% =============================================================
\thispagestyle{empty}
\begin{center}
{\bfseries UCAPAN TERIMA KASIH\par}
\vspace{10mm}
\end{center}

\noindent Puji syukur penulis panjatkan kepada Tuhan Yang Maha Esa atas rahmat dan karunia-Nya sehingga tesis yang berjudul \textit{``Judul Tesis untuk Program Magister (S2) Data Science dan Kecerdasan Buatan''} ini dapat diselesaikan sebagai salah satu syarat memperoleh gelar Magister pada Program Studi S2 Data Science dan Kecerdasan Buatan, Fakultas Ilmu Komputer dan Teknologi Informasi, Universitas Sumatera Utara.

Penulis mengucapkan terima kasih kepada Komisi Pembimbing atas bimbingan, arahan, dan waktu yang telah diberikan selama proses penelitian dan penulisan tesis ini. Ucapan terima kasih juga disampaikan kepada seluruh dosen dan staf Program Studi S2 Data Science dan Kecerdasan Buatan, keluarga, serta rekan-rekan yang telah memberikan dukungan moril maupun materiil.

Penulis menyadari bahwa tesis ini masih memiliki kekurangan. Oleh karena itu, kritik dan saran yang membangun sangat diharapkan demi kesempurnaan penelitian selanjutnya.

\vspace{10mm}
\noindent Medan, \the\year

\vspace{6mm}
\noindent Penulis
\clearpage

% =============================================================
% Abstrak, based on Lampiran 10 of the Pedoman.
% One paragraph only, 200-300 words: latar belakang, tujuan,
% metodologi, dan hasil. No tables/figures/direct quotations.
% =============================================================
\thispagestyle{empty}
\begin{center}
{\bfseries ABSTRAK\par}
\vspace{10mm}
\end{center}

\noindent Perkembangan volume data yang pesat pada berbagai sektor menuntut pendekatan analisis yang mampu mengekstraksi pola dan pengetahuan secara otomatis dan akurat. Penelitian ini dilatarbelakangi oleh permasalahan tersebut, khususnya pada domain \textit{[nama domain aplikasi, misalnya: deteksi anomali/klasifikasi teks/prediksi deret waktu]}, di mana pendekatan konvensional belum mampu memberikan hasil yang optimal dari segi akurasi maupun efisiensi komputasi. Penelitian ini bertujuan untuk merancang dan mengevaluasi sebuah pendekatan berbasis \textit{[nama metode/algoritma, misalnya: pembelajaran mendalam, model probabilistik Bayesian, atau algoritma metaheuristik]} yang dapat meningkatkan kinerja pada permasalahan tersebut dibandingkan dengan metode \textit{baseline}. Metodologi penelitian terdiri atas tahap pengumpulan dan praproses data, perancangan arsitektur model, pelatihan dan penyetelan hiperparameter, serta pengujian menggunakan skema validasi silang pada himpunan data \textit{[nama/asal dataset]}. Kinerja model dievaluasi menggunakan metrik \textit{[nama metrik, misalnya: akurasi, F1-score, RMSE]} dan dibandingkan dengan beberapa metode pembanding dari literatur terkini. Hasil penelitian menunjukkan bahwa metode yang diusulkan memberikan peningkatan kinerja sebesar \textit{[nilai/persentase]} dibandingkan dengan metode pembanding, serta menunjukkan konsistensi hasil pada beberapa skenario pengujian. Temuan ini mengindikasikan bahwa pendekatan yang diusulkan berpotensi diterapkan pada permasalahan sejenis dalam skala yang lebih luas, sekaligus membuka peluang penelitian lanjutan terkait optimalisasi arsitektur model dan efisiensi komputasi.

\vspace{6mm}
\noindent\textbf{Kata Kunci:} kata kunci satu, kata kunci dua, kata kunci tiga, kata kunci empat.
\clearpage

% =============================================================
% Abstract (English translation of Abstrak), Lampiran 11.
% Title = translated thesis title, then the word ABSTRACT, then
% one paragraph, 200-300 words.
% =============================================================
\thispagestyle{empty}
\begin{center}
\textit{Thesis Title in English}\par
\vspace{6mm}
{\bfseries ABSTRACT\par}
\vspace{10mm}
\end{center}

\noindent The rapid growth of data volume across various sectors demands analytical approaches capable of automatically and accurately extracting patterns and knowledge. This research is motivated by that challenge, specifically within the domain of \textit{[application domain, e.g., anomaly detection/text classification/time-series forecasting]}, where conventional approaches have not been able to deliver optimal results in terms of either accuracy or computational efficiency. This study aims to design and evaluate an approach based on \textit{[method/algorithm name, e.g., deep learning, a Bayesian probabilistic model, or a metaheuristic algorithm]} that improves performance on the problem compared to a baseline method. The research methodology consists of data collection and preprocessing, model architecture design, training and hyperparameter tuning, and evaluation using a cross-validation scheme on the \textit{[dataset name/source]} dataset. Model performance is evaluated using \textit{[metric name, e.g., accuracy, F1-score, RMSE]} and compared against several state-of-the-art baseline methods. The results show that the proposed method achieves a performance improvement of \textit{[value/percentage]} over the baseline methods, with consistent results across several testing scenarios. These findings indicate that the proposed approach has the potential to be applied to similar problems at a larger scale, while also opening opportunities for further research on model architecture optimization and computational efficiency.

\vspace{6mm}
\noindent\textbf{Keywords:} keyword one, keyword two, keyword three, keyword four.
\clearpage


\tableofcontents
\listoftables
\listoffigures

% ---------------------------------------------------------------
% BAGIAN TUBUH (body) -- arabic page numbers from BAB 1 onward.
% ---------------------------------------------------------------
\pagenumbering{arabic}
\chapter{PENDAHULUAN}

\section{Latar Belakang}
Bab Pendahuluan adalah bab terpendek dalam sebuah tesis dan merupakan bab yang paling penting. Satu-satunya tujuan penulisan bab ini adalah untuk memperkenalkan kepada pembaca tentang apa yang akan dipaparkan dalam tesis \citep{russell2020}.

Jelaskan lebih dahulu apa yang memotivasi penelitian yang akan dilakukan dan mengapa penelitian pada domain \textit{[nama domain aplikasi]} tersebut perlu dilakukan. Paparkan secara sistematis gejala dan peristiwa yang menimbulkan permasalahan untuk diteliti, serta sebutkan penelitian-penelitian terdahulu yang relevan pada bidang yang akan diteliti \citep{goh2009}.

\subsection{Ruang lingkup permasalahan yang diteliti}
Pada bagian ini penulis menjelaskan secara ringkas gambaran umum permasalahan penelitian, batasan awal ruang lingkupnya, serta perbedaan kontribusi penelitian ini dibandingkan dengan penelitian-penelitian sebelumnya yang telah disebutkan.

\section{Rumusan Masalah}
Rumusan masalah dapat ditulis dalam bentuk \textit{research question} atau \textit{problem statement}. Contoh rumusan masalah berbentuk pertanyaan penelitian:
\begin{enumerate}
    \item Bagaimana pengaruh \textit{[variabel/metode]} terhadap kinerja \textit{[tugas/permasalahan]} pada \textit{[domain aplikasi]}?
    \item Bagaimana perbandingan kinerja pendekatan yang diusulkan terhadap metode \textit{baseline} yang ada?
\end{enumerate}

\section{Tujuan Penelitian}
Tujuan penelitian ini adalah untuk merancang, mengimplementasikan, dan mengevaluasi \textit{[nama metode/algoritma]} guna menyelesaikan permasalahan yang telah dirumuskan pada Subbab 1.2.

\section{Manfaat Penelitian}
Manfaat yang diharapkan dari penelitian ini antara lain memberikan kontribusi teoritis pada bidang \textit{[bidang keilmuan]} serta manfaat praktis berupa \textit{[manfaat terapan, misalnya: purwarupa sistem, rekomendasi kebijakan, dsb.]}.

\section{Batasan Masalah}
Agar penelitian tetap fokus, ruang lingkup penelitian dibatasi pada hal-hal berikut:
\begin{enumerate}
    \item Penelitian dibatasi pada disiplin ilmu data science dan kecerdasan buatan.
    \item Data yang digunakan adalah \textit{[nama/sumber dataset]}.
    \item Evaluasi dilakukan menggunakan metrik \textit{[nama metrik]} dan tidak mencakup \textit{[hal yang dikecualikan]}.
\end{enumerate}

\section{Sistematika Penulisan}
Sistematika penulisan tesis ini terdiri atas lima bab. Bab 1 berisi pendahuluan. Bab 2 membahas tinjauan pustaka dan landasan teori. Bab 3 menguraikan metodologi penelitian. Bab 4 menyajikan hasil dan pembahasan. Bab 5 berisi kesimpulan dan saran.

\chapter{TINJAUAN PUSTAKA}

\section{Penelitian Terdahulu}
Beberapa penelitian sebelumnya telah mengeksplorasi permasalahan serupa. \citet{riedmiller1994} mengusulkan algoritma \textit{gradient descent} untuk meminimumkan \textit{mean square error} (MSE) pada pelatihan jaringan saraf tiruan. Pendekatan ini kemudian dikembangkan lebih lanjut oleh beberapa peneliti lain \citep{calvert2005,aaron2014}.

Sementara itu, meskipun dianggap sebagai pendekatan yang relatif baru, algoritma \textit{artificial bee colony} (ABC) telah digunakan secara luas dalam pencarian penyelesaian optimum \citep{shukran2011}.

\subsection{Ringkasan perbandingan penelitian terdahulu}
Tabel~\ref{tab:perbandingan} merangkum perbandingan pendekatan, dataset, dan metrik evaluasi yang digunakan pada beberapa penelitian terdahulu yang relevan dengan penelitian ini.

\begin{table}[H]
\centering
\caption{Perbandingan penelitian terdahulu}
\label{tab:perbandingan}
\begin{tabular}{@{}lllc@{}}
\toprule
Peneliti & Metode & Dataset & Metrik \\
\midrule
\citet{riedmiller1994} & Gradient Descent & Sintetis & MSE \\
\citet{shukran2011}    & Artificial Bee Colony & Benchmark & Akurasi \\
\citet{russell2020}    & Pembelajaran Mesin & Umum & F1-score \\
\bottomrule
\end{tabular}
\end{table}

\section{Landasan Teori}

\subsection{Data science dan kecerdasan buatan}
Data science merupakan bidang interdisipliner yang menggabungkan statistik, ilmu komputer, dan pengetahuan domain untuk mengekstraksi wawasan dari data. Kecerdasan buatan (\textit{artificial intelligence}) adalah cabang ilmu komputer yang berfokus pada pembangunan sistem yang mampu meniru kemampuan kognitif manusia \citep{russell2020}.

\subsection{Metode yang relevan dengan penelitian ini}
Uraikan di sini teori dan konsep matematis dari metode/algoritma yang menjadi dasar pendekatan yang diusulkan pada Bab 3, misalnya arsitektur model pembelajaran mesin, fungsi \textit{loss}, atau kerangka evaluasi yang digunakan.

\chapter{METODOLOGI PENELITIAN}

\section{Tahapan Penelitian}
Penelitian ini dilaksanakan melalui beberapa tahapan utama: (1) pengumpulan dan praproses data, (2) perancangan model, (3) pelatihan dan penyetelan hiperparameter, serta (4) pengujian dan evaluasi.

\subsection{Kerangka kerja penelitian}
Kerangka kerja penelitian secara umum mengikuti siklus iteratif antara perancangan model dan evaluasi hasil, hingga diperoleh konfigurasi model dengan kinerja terbaik pada data validasi.

\section{Pengumpulan dan Praproses Data}
Data yang digunakan dalam penelitian ini bersumber dari \textit{[nama/sumber dataset]}. Tahap praproses meliputi pembersihan data, penanganan nilai hilang, normalisasi fitur, dan pembagian data menjadi himpunan latih, validasi, dan uji.

\section{Perancangan Model}
Diberikan sebuah himpunan data $D = \{(x_i, y_i)\}_{i=1}^{n}$, model yang diusulkan bertujuan mempelajari fungsi $f_\theta: \mathcal{X} \rightarrow \mathcal{Y}$ yang meminimumkan fungsi kerugian (\textit{loss function}) berikut:
\begin{equation}
\label{eq:loss}
\mathcal{L}(\theta) = \frac{1}{n}\sum_{i=1}^{n} \ell\big(f_\theta(x_i), y_i\big) + \lambda \lVert \theta \rVert_2^2
\end{equation}
dengan $\ell(\cdot,\cdot)$ adalah fungsi kerugian dasar, $\theta$ adalah parameter model, dan $\lambda$ adalah koefisien regularisasi.

Parameter optimal $\theta^{*}$ diperoleh melalui optimisasi berbasis gradien sebagaimana pada Persamaan~\eqref{eq:update}:
\begin{align}
\label{eq:update}
\theta^{(t+1)} &= \theta^{(t)} - \eta \, \nabla_\theta \mathcal{L}\big(\theta^{(t)}\big)
\end{align}
dengan $\eta$ adalah laju pembelajaran (\textit{learning rate}) dan $t$ adalah indeks iterasi.

\subsection{Arsitektur model yang diusulkan}
Jelaskan di sini arsitektur model secara rinci, termasuk lapisan/komponen penyusun, fungsi aktivasi yang digunakan, serta alasan pemilihan arsitektur tersebut dikaitkan dengan tinjauan pustaka pada Bab 2.

\section{Skenario Pengujian}
Evaluasi kinerja model dilakukan menggunakan skema validasi silang \textit{k-fold} dengan metrik evaluasi berupa \textit{[nama metrik]}. Hasil pengujian dibandingkan dengan beberapa metode \textit{baseline} sebagaimana diuraikan pada Bab 2.

\chapter{HASIL DAN PEMBAHASAN}

\section{Karakteristik Data}
Bagian ini memaparkan karakteristik data hasil praproses yang digunakan pada tahap pelatihan dan pengujian model, sebagaimana diuraikan pada Subbab 3.2.

\begin{table}[H]
\centering
\caption{Statistik ringkas himpunan data}
\label{tab:statistik-data}
\begin{tabular}{@{}lccc@{}}
\toprule
Himpunan & Jumlah Sampel & Jumlah Fitur & Proporsi (\%) \\
\midrule
Latih    & 0 & 0 & 70 \\
Validasi & 0 & 0 & 15 \\
Uji      & 0 & 0 & 15 \\
\bottomrule
\end{tabular}
\end{table}

\section{Hasil Pelatihan Model}
Gambar~\ref{fig:kurva-pelatihan} memperlihatkan kurva \textit{loss} pelatihan dan validasi terhadap jumlah iterasi (\textit{epoch}). Kurva ini digunakan untuk memantau kekonvergenan model serta mendeteksi indikasi \textit{overfitting}.

\begin{figure}[H]
\centering
\begin{tikzpicture}
  \draw[->] (0,0) -- (7,0) node[right] {\footnotesize Epoch};
  \draw[->] (0,0) -- (0,4) node[above] {\footnotesize Loss};
  \draw[thick] (0,3.5) .. controls (2,1.2) and (4,0.6) .. (6.5,0.4);
  \draw[thick,dashed] (0,3.8) .. controls (2,1.6) and (4,1.0) .. (6.5,0.9);
  \node[anchor=west] at (5,1.6) {\footnotesize \textit{training loss}};
  \node[anchor=west] at (5,2.4) {\footnotesize \textit{validation loss}};
\end{tikzpicture}
\caption{Kurva \textit{loss} pelatihan dan validasi (ilustrasi)}
\label{fig:kurva-pelatihan}
\end{figure}

\section{Perbandingan Kinerja dengan Metode Pembanding}
Tabel~\ref{tab:hasil-perbandingan} menyajikan perbandingan kinerja pendekatan yang diusulkan dengan beberapa metode \textit{baseline} yang telah disebutkan pada Bab 2 \citep{russell2020,shukran2011}.

\begin{table}[H]
\centering
\caption{Perbandingan kinerja model}
\label{tab:hasil-perbandingan}
\begin{tabular}{@{}lccc@{}}
\toprule
Metode & Akurasi & F1-score & Waktu Komputasi (s) \\
\midrule
\textit{Baseline} 1 & 0.00 & 0.00 & 0.00 \\
\textit{Baseline} 2 & 0.00 & 0.00 & 0.00 \\
Metode yang Diusulkan & \textbf{0.00} & \textbf{0.00} & 0.00 \\
\bottomrule
\end{tabular}
\end{table}

\section{Pembahasan}
Hasil pada Tabel~\ref{tab:hasil-perbandingan} menunjukkan bahwa metode yang diusulkan memberikan kinerja \textit{[lebih baik/setara/dst.]} dibandingkan metode pembanding. Diskusikan di sini keterkaitan hasil dengan rumusan masalah pada Subbab 1.2 dan tinjauan pustaka pada Bab 2, termasuk kemungkinan penjelasan atas hasil yang diperoleh serta keterbatasannya.

\chapter{KESIMPULAN DAN SARAN}

\section{Kesimpulan}
Berdasarkan hasil penelitian dan pembahasan pada Bab 4, dapat disimpulkan bahwa:
\begin{enumerate}
    \item Metode yang diusulkan pada penelitian ini mampu menjawab tujuan penelitian sebagaimana dirumuskan pada Subbab 1.3.
    \item Pendekatan yang diusulkan menunjukkan kinerja \textit{[lebih baik/setara]} dibandingkan metode pembanding pada metrik evaluasi yang digunakan.
\end{enumerate}

\section{Saran}
Beberapa saran untuk penelitian selanjutnya antara lain:
\begin{enumerate}
    \item Mengeksplorasi \textit{[arsitektur/metode lain]} untuk meningkatkan kinerja lebih lanjut.
    \item Menguji generalisasi pendekatan yang diusulkan pada \textit{dataset} lain di luar cakupan penelitian ini.
    \item Mengkaji efisiensi komputasi model untuk kebutuhan penerapan pada skala produksi.
\end{enumerate}


% ---------------------------------------------------------------
% BAGIAN BELAKANG (back matter) -- numbering continues from the body.
% ---------------------------------------------------------------
{\singlespacing
\bibliography{references}
}

\appendix
% Appendix chapters are labelled "LAMPIRAN A", "LAMPIRAN B", ... instead of "BAB".
\titleformat{\chapter}[display]
  {\normalfont\bfseries\centering}
  {\Large LAMPIRAN \thechapter}
  {12pt}
  {\Large\MakeUppercase}
\chapter{Kode Program}

Lampiran ini memuat \textit{source code} yang dibangun sebagai hasil implementasi model yang dikemukakan pada Bab 3.

\begin{lstlisting}[language=Python, basicstyle=\ttfamily\small, breaklines=true, frame=single, caption={Contoh cuplikan kode pelatihan model}]
def train(model, dataloader, optimizer, loss_fn, epochs):
    for epoch in range(epochs):
        for x_batch, y_batch in dataloader:
            optimizer.zero_grad()
            y_pred = model(x_batch)
            loss = loss_fn(y_pred, y_batch)
            loss.backward()
            optimizer.step()
    return model
\end{lstlisting}

\chapter{Data Tambahan}

Lampiran ini memuat data mentah, tabel hasil pengujian lengkap, atau materi pendukung lain yang tidak dimuat pada teks utama, sebagaimana diuraikan pada Subbab 6.4 Pedoman Penulisan Tesis.

\begin{table}[H]
\centering
\caption{Hasil pengujian lengkap per fold validasi silang}
\label{tab:lampiran-hasil-lengkap}
\begin{tabular}{@{}cccc@{}}
\toprule
Fold & Akurasi & F1-score & Waktu (s) \\
\midrule
1 & 0.00 & 0.00 & 0.00 \\
2 & 0.00 & 0.00 & 0.00 \\
3 & 0.00 & 0.00 & 0.00 \\
4 & 0.00 & 0.00 & 0.00 \\
5 & 0.00 & 0.00 & 0.00 \\
\bottomrule
\end{tabular}
\end{table}


\end{document}
